%-*-tex-*-
% Copyright Michael J. Ferguson, INRS-Telecommunications
% All rights reserved. 

% ======== Master Multilingual Pattern Inputs ======

% This file defines the changes in the \lccodes needed for Multilingual TeX
% and then inputs the patterns for English and French. 

% =========== Special lccodes and hyphenation for Bilingual Version ====


\gdef\accenthyphcodes{
\def\oe{^^[} % \oe 
\def\i{^^P}
\def\'##1{\csname @ac@##1\endcsname}
\def\`##1{\csname @gr@##1\endcsname}
\def\v##1{\csname @v@##1\endcsname}
\let\^^_=\v
\def\u##1{\csname  @u@##1\endcsname}
\let\^^S=\u
\def\=##1{\csname @eq@##1\endcsname}
\def\^##1{\csname @hat@##1\endcsname}
\let\^^D=\^
\def\.##1{\csname @dot@##1\endcsname}
\def\H##1{\csname @H@##1\endcsname}
\def\~##1{\csname @til@##1\endcsname}
\def\"##1{\csname @um@##1\endcsname}
\let\c@@=\c
\def\c##1{\csname c@##1\endcsname}
}

\gdef\spechyphcodes{}



% ======== Pattern Input / English and French ======

% english hyphenation patterns
\begingroup  
\language=0
\input enhyph \relax
\endgroup

% french hyphenation patterns
\begingroup
\language=1
\input frhyph \relax
\endgroup

% english hyphenation exceptions
\begingroup
\language=0
\input enhyphex \relax
\endgroup

% ======= definitions for fast eng/fr hyphenation =====
\def\ehyph{\language=0 \lccode`\'=0 \nonfrenchspacing}
\def\fhyph{\language=1 \lccode`\'=`\'\frenchspacing}